\subsection{Number fields and arbitrary nonnegative shifts}\label{sec:dedekind-shift}
Let $K$ be a number field, let $a_K(n)$ count its nonzero integral ideals of norm $n$, and put
\[
 Z_{K,t}(s)=\sum_{n\ge1}a_K(n)(n+t)^{-s},\qquad t\ge0,\quad\Re s>1.
\]
Thus $Z_{K,0}=\zeta_K$. We use the classical continuation of $\zeta_K$, whose sole pole is simple at $1$, with positive residue $\kappa_K$; see \citet[Chapter VII, \S5, Proposition 5.2 and Corollary 5.11]{NeukirchShift}.

\begin{proposition}[The shifted ideal series]\label{prop:dedekind-all-shifts}
For every fixed $t\ge0$, the difference $Z_{K,t}-\zeta_K$ extends to an entire function of $s$. Consequently $Z_{K,t}$ has exactly one pole, at $s=1$, with residue $\kappa_K$. The series of finite congruence quotients of $(\mathcal O_K)_0$, counted by their kernel congruences and weighted by their cardinalities, is
\[
 \mathcal Q_K(s)=\zeta_K(s)+Z_{K,1}(s),
\]
and its sole pole is simple at $1$, with residue $2\kappa_K$.
If $t\ge0$ is not an integer, $Z_{K,t}$ cannot equal an ordinary integer-base Dirichlet series on any half-plane of absolute convergence.
\end{proposition}
\begin{proof}
Choose an integer $A>\max(1,t)$ and remove the finitely many small norms:
\[
 H_A(s)=\zeta_K(s)-\sum_{n<A}a_K(n)n^{-s}.
\]
The binomial expansion gives, initially for $\Re s>1$,
\begin{equation}\label{eq:dedekind-cutoff}
 Z_{K,t}(s)=\sum_{n<A}a_K(n)(n+t)^{-s}
 +\sum_{m\ge0}\frac{(-1)^m(s)_m}{m!}t^m H_A(s+m).
\end{equation}
Absolute convergence permits interchanging the two sums, since $t/n\le t/A<1$ in the tail.

For $m\ge1$, the only possible pole of $H_A(s+m)$ is at $s=1-m$. One of the factors of $(s)_m=s(s+1)\cdots(s+m-1)$ vanishes there, so each such summand is entire. We also need convergence after this cancellation. Fix a compact set of $s$-values and choose an integer $R\ge1$ with $|s|\le R$ on it. For $m\ge R+2$,
\[
 |H_A(s+m)|\le\sum_{n\ge A}a_K(n)n^{-m+R}
 \le A^{R+2-m}\zeta_K(2),
\]
while
\[
 \frac{|(s)_m|}{m!}\le\frac{(R)_m}{m!}
 =\binom{m+R-1}{R-1}.
\]
The large-$m$ terms are therefore bounded by a fixed polynomial in $m$ times $(t/A)^m$. This proves normal convergence on every compact $s$-set; the finitely many remaining terms are entire after the displayed cancellations. Subtracting $\zeta_K(s)$ from the $m=0$ term and the finite first sum in \eqref{eq:dedekind-cutoff} leaves an entire function. This proves the continuation and residue statements, including the endpoint $t=1$.

Lemma~\ref{lem:congruences} gives exactly two finite congruences for each nonzero ideal $J\subseteq\mathcal O_K$, with quotient cardinalities $N(J)$ and $N(J)+1$. Summing them proves the formula for $\mathcal Q_K$ and doubles the residue. This counts congruences, so different ideals with isomorphic quotient rings are still counted separately.

Finally $a_K(1)=1$, and hence
\[
 \lim_{\sigma\to+\infty}(1+t)^\sigma Z_{K,t}(\sigma)=1.
\]
The limit follows by dominated convergence after fixing one absolute-convergence abscissa. A nonzero ordinary Dirichlet series has a least integer index $n_0$ with nonzero coefficient, and the same argument identifies its leading exponential scale as $n_0$. Equality with $Z_{K,t}$ would force $n_0=1+t$, which is impossible when $t$ is not an integer.
\end{proof}

For $0\le t<1$ one may take the cutoff before the first norm, recovering the simpler expansion
\[
 Z_{K,t}(s)=\sum_{m\ge0}\frac{(-1)^m(s)_m}{m!}t^m\zeta_K(s+m).
\]
The cutoff in \eqref{eq:dedekind-cutoff} supplies the convergence needed at larger shifts. For any $t>0$ the shifted count tends to zero as $s\to+\infty$ on the real axis, whereas a normalized absolutely convergent Euler product tends to one. Thus the shifted cardinality convention does not have such an Euler product.
